% GENERATED FILE. Edit canonical lessons, then run npm run generate:books.
% canonical-source-hash: fnv1a64:d73f82baeaf627c5
% canonical-lessons: PA-C145-bahadar, PA-C145-shat, PA-C145-bhukkha, PA-C145-piasa, PA-C145-taza

\chapter{Brave, Calm, Hungry, Thirsty, Fresh}
\label{ch:pa-brave-calm-hungry-thirsty-fresh}
\hlchaptermodality{145}

\begin{chapteropening}
\textbf{By the end of this chapter:} \emph{I can say brave, calm, hungry, thirsty and fresh.}

Complete the last lesson of chapter 145: i can say brave, calm, hungry, thirsty and fresh.
\end{chapteropening}

\section[\texorpdfstring{\pa{ਬਹਾਦਰ} (bahādar)}{ਬਹਾਦਰ (bahādar)}]{\pa{ਬਹਾਦਰ} (bahādar) — brave}

\label{lesson:PA-C145-bahadar}



\begin{quote}
Before the new one: say the Punjabi for clever, then the Punjabi for lazy.
\end{quote}

\subsection*{You'll want to know: \pa{ਬਹਾਦਰ}}
\textbf{\pa{ਬਹਾਦਰ}} — \emph{bahādar} — \textquotedblleft{}brave\textquotedblright{}.

\begin{grammarlens}[title={What you've built}]
One more word for this chapter.
\end{grammarlens}

\subsection*{Your turn}
\begin{itemize}
\raggedright
  \item \emph{Say it:} \emph{bahādar}
  \item \emph{Say it:} \emph{bahādar}, once more
  \item \emph{Say it:} say \emph{bahādar}
  \item \emph{Recall:} say the Punjabi for clever, then the Punjabi for lazy, then say \emph{bahādar} again
\end{itemize}

\begin{tcolorbox}[breakable,colback=teal!4,colframe=teal!35!black,title={Before you move on}]
What does \pa{ਬਹਾਦਰ} mean? (\textquotedblleft{}Brave\textquotedblright{}.) Say it once more.
\end{tcolorbox}
\section[\texorpdfstring{\pa{ਸ਼ਾਂਤ} (shā̃t)}{ਸ਼ਾਂਤ (shā̃t)}]{\pa{ਸ਼ਾਂਤ} (shā̃t) — calm, quiet}

\label{lesson:PA-C145-shat}



\begin{quote}
Before the new one: say the Punjabi for lazy, then the Punjabi for brave.
\end{quote}

\subsection*{You'll want to know: \pa{ਸ਼ਾਂਤ}}
\textbf{\pa{ਸ਼ਾਂਤ}} — \emph{shā̃t} — \textquotedblleft{}calm, quiet\textquotedblright{}.

\begin{grammarlens}[title={What you've built}]
One more word for this chapter.
\end{grammarlens}

\subsection*{Your turn}
\begin{itemize}
\raggedright
  \item \emph{Say it:} \emph{shā̃t}
  \item \emph{Say it:} \emph{shā̃t}, once more
  \item \emph{Say it:} say \emph{shā̃t}
  \item \emph{Recall:} say the Punjabi for lazy, then the Punjabi for brave, then say \emph{shā̃t} again
\end{itemize}

\begin{tcolorbox}[breakable,colback=teal!4,colframe=teal!35!black,title={Before you move on}]
What does \pa{ਸ਼ਾਂਤ} mean? (\textquotedblleft{}Calm, quiet\textquotedblright{}.) Say it once more.
\end{tcolorbox}
\section[\texorpdfstring{\pa{ਭੁੱਖਾ} (bhukkhā)}{ਭੁੱਖਾ (bhukkhā)}]{\pa{ਭੁੱਖਾ} (bhukkhā) — hungry}

\label{lesson:PA-C145-bhukkha}



\begin{quote}
Before the new one: say the Punjabi for brave, then the Punjabi for calm.
\end{quote}

\subsection*{You'll want to know: \pa{ਭੁੱਖਾ}}
\textbf{\pa{ਭੁੱਖਾ}} — \emph{bhukkhā} — \textquotedblleft{}hungry\textquotedblright{}.

\begin{grammarlens}[title={What you've built}]
One more word for this chapter.
\end{grammarlens}

\subsection*{Your turn}
\begin{itemize}
\raggedright
  \item \emph{Say it:} \emph{bhukkhā}
  \item \emph{Say it:} \emph{bhukkhā}, once more
  \item \emph{Say it:} say \emph{bhukkhā}
  \item \emph{Recall:} say the Punjabi for brave, then the Punjabi for calm, then say \emph{bhukkhā} again
\end{itemize}

\begin{tcolorbox}[breakable,colback=teal!4,colframe=teal!35!black,title={Before you move on}]
What does \pa{ਭੁੱਖਾ} mean? (\textquotedblleft{}Hungry\textquotedblright{}.) Say it once more.
\end{tcolorbox}
\section[\texorpdfstring{\pa{ਪਿਆਸਾ} (piāsā)}{ਪਿਆਸਾ (piāsā)}]{\pa{ਪਿਆਸਾ} (piāsā) — thirsty}

\label{lesson:PA-C145-piasa}



\begin{quote}
Before the new one: say the Punjabi for calm, then the Punjabi for hungry.
\end{quote}

\subsection*{You'll want to know: \pa{ਪਿਆਸਾ}}
\textbf{\pa{ਪਿਆਸਾ}} — \emph{piāsā} — \textquotedblleft{}thirsty\textquotedblright{}.

\begin{grammarlens}[title={What you've built}]
One more word for this chapter.
\end{grammarlens}

\subsection*{Your turn}
\begin{itemize}
\raggedright
  \item \emph{Say it:} \emph{piāsā}
  \item \emph{Say it:} \emph{piāsā}, once more
  \item \emph{Say it:} say \emph{piāsā}
  \item \emph{Recall:} say the Punjabi for calm, then the Punjabi for hungry, then say \emph{piāsā} again
\end{itemize}

\begin{tcolorbox}[breakable,colback=teal!4,colframe=teal!35!black,title={Before you move on}]
What does \pa{ਪਿਆਸਾ} mean? (\textquotedblleft{}Thirsty\textquotedblright{}.) Say it once more.
\end{tcolorbox}
\section[\texorpdfstring{\pa{ਤਾਜ਼ਾ} (tāzā)}{ਤਾਜ਼ਾ (tāzā)}]{\pa{ਤਾਜ਼ਾ} (tāzā) — fresh}

\label{lesson:PA-C145-taza}



\begin{quote}
Before the new one: say the Punjabi for hungry, then the Punjabi for thirsty.
\end{quote}

\subsection*{You'll want to know: \pa{ਤਾਜ਼ਾ}}
\textbf{\pa{ਤਾਜ਼ਾ}} — \emph{tāzā} — \textquotedblleft{}fresh\textquotedblright{}.

\begin{grammarlens}[title={What you've built}]
One more word for this chapter.
\end{grammarlens}

\subsection*{Your turn}
\begin{itemize}
\raggedright
  \item \emph{Say it:} \emph{tāzā}
  \item \emph{Say it:} \emph{tāzā}, once more
  \item \emph{Say it:} say \emph{tāzā}
  \item \emph{Recall:} say the Punjabi for hungry, then the Punjabi for thirsty, then say \emph{tāzā} again
\end{itemize}

\begin{tcolorbox}[breakable,colback=teal!4,colframe=teal!35!black,title={Before you move on}]
What does \pa{ਤਾਜ਼ਾ} mean? (\textquotedblleft{}Fresh\textquotedblright{}.) Say it once more.
\end{tcolorbox}
